%% =====================================================================
%%  structures/oficio/structure.tex  --  Esqueleto retorico del OFICIO
%% ---------------------------------------------------------------------
%%  Macros que componen las partes del oficio en el ORDEN del machote:
%%
%%     \lugaryfecha                       Lima, 30 de julio del 2026
%%     \numerodocumento                   OFICIO N° 02-2026/ABC/CR
%%     \destinatario{Señor:}{Nombre}{cargo-o-vacio}
%%     \asunto{...}   \referencia{...}
%%     \saludo                            De mi especial consideracion:
%%     ... parrafos libres ...
%%     \atentamente   \bloquefirma
%%     \adjunto{...}  \concopia{...}
%%
%%  Los datos (numero, anio, fecha, lugar) se declaran con \setmeta en
%%  el metadata.tex del documento; aqui solo se COMPONEN.
%%
%%  OJO: este archivo se carga en el preambulo del documento; no usar
%%  nombres con @ ni \newcommand duplicables (guardia anti-doble-carga).
%% =====================================================================
\ifcsname despachoEstructuraOficio\endcsname
  \expandafter\endinput
\fi
\expandafter\def\csname despachoEstructuraOficio\endcsname{cargada}

%% --- Ajustes finos (sobrescribibles en el documento) -----------------
\newlength{\sangriaasunto}
%% Desde la MITAD de la hoja hacia la derecha, como los oficios de la bancada (Edison
%% 2026-10-10); antes 3,20 cm. La etiqueta a 5,8 cm: el texto, tras «REFERENCIA: », arranca en la
%% mitad de la caja (15,92 cm en A4 con 2,54 de margen).
\setlength{\sangriaasunto}{7.96cm}   % sangria del bloque ASUNTO/REF: el TEXTO arranca en la mitad de la hoja
                                     % (10,5 cm del borde = 7,96 cm del margen); la etiqueta cuelga a su izquierda
                                     % (Edison 2026-10-10, segunda vuelta: «el texto es lo que comienza del medio»)
\newcommand{\setsangriaasunto}[1]{\setlength{\sangriaasunto}{#1}}

%% --- Aire vertical ENTRE las partes de la cabecera -------------------
%%  UN solo mando para el espacio que separa fecha, numero, destinatario
%%  y asunto/referencia. Son FACTORES (no longitudes) para que se
%%  evaluen en el punto de uso y sigan al interlineado del documento:
%%  una \newlength congelaria el \baselineskip del preambulo y
%%  \setinterlineado dejaria de afectarlos.
%%  OJO: al aire de cada separacion se le SUMA el \parskip del documento
%%  (6pt por defecto, ver core/config/espaciado.tex), asi que el hueco
%%  real es mayor que el factor.
%%  Bajados de 0.8-0.9 a 0.20 y el saludo de 0.4 a 0.10 -- pedido de
%%  Edison 2026-09-07: los blancos de la cabecera «ocupan mucho» y habia
%%  que reducirlos en la configuracion del sistema, no documento a
%%  documento (primera pasada 0.35/0.20/0.25; devuelto con «reducir mas»,
%%  junto con el margen superior y el \parskip). Sobrescribibles en el
%%  preambulo de un .tex a mano.
\newcommand{\espaciopartes}{0.20}
\newcommand{\setespaciopartes}[1]{\renewcommand{\espaciopartes}{#1}}
\newcommand{\espaciosaludo}{0.10}
\newcommand{\setespaciosaludo}[1]{\renewcommand{\espaciosaludo}{#1}}
%%  Aire interno del bloque destinatario, antes de «Con copia a:»
%%  (dentro de la minipage el \parskip es 2pt, no 6pt).
\newcommand{\espaciocopiadest}{0.15}
\newcommand{\setespaciocopiadest}[1]{\renewcommand{\espaciocopiadest}{#1}}

%% --- Lugar y fecha (derecha) -----------------------------------------
%%  Usa \meta{fecha} si existe; si no, la fecha de compilacion.
\newcommand{\lugaryfecha}{%
  \noindent
  {\raggedleft
   \meta{lugar}, \ifmetadef{fecha}{\meta{fecha}}{\today}\par}%
  \vspace{\espaciopartes\baselineskip}%
}

%% --- Numero de documento (subrayado, como el machote) ----------------
%%  Si \setmeta{numerocompleto}{...} esta definido, se usa TAL CUAL
%%  (numeraciones no estandar, p. ej. oficios multiples); si no, se
%%  compone tipodoc + numero-anio/siglas.
\newcommand{\numerodocumento}{%
  \noindent
  \ifmetadef{numerocompleto}%
    {\underline{\MakeUppercase{\meta{numerocompleto}}}}%
    {\underline{\MakeUppercase{\meta{tipodoc}}\ N°~\meta{numero}-\meta{anio}/\meta{siglas}}}%
  \par\vspace{\espaciopartes\baselineskip}%
}

%% --- Destinatario -----------------------------------------------------
%%  \destinatario{Señor:}{NOMBRE APELLIDOS}{Cargo (o vacio)}
%%  Compone tambien el "Presente. –" final.
%%  \destinatario[bloque-copia]{Señor:}{NOMBRE}{Cargo}
%%  El argumento OPCIONAL es el bloque «Con copia a:» que va arriba (antes
%%  de «Presente.»); si se omite, el destinatario es el normal. El tamaño
%%  del bloque se ajusta con \settamdestinatario (\despachotamdest).
%% Institucion o entidad del destinatario: CUARTA linea del bloque, debajo
%% del cargo (p. ej. «Taller Cultural de la UNE - La Cantuta (TACUNE)»). Se
%% fija aparte para no cambiar la aridad de \destinatario y que los .tex ya
%% compilados sigan valiendo (2026-09-23). Vacia por defecto: no imprime nada.
\newcommand{\destinatarioinstitucion}{}
\newcommand{\setdestinatarioinstitucion}[1]{\renewcommand{\destinatarioinstitucion}{#1}}

\newcommand{\destinatario}[4][]{%
  \noindent
  \begin{minipage}{\textwidth}%
    \setlength{\parskip}{2pt}%
    \despachotamdest\despachointerdest
    #2\par
    {\seriesemi #3}\par
    \if\relax\detokenize{#4}\relax\else
      #4\par
    \fi
    \ifx\destinatarioinstitucion\empty\else
      \destinatarioinstitucion\par
    \fi
    \if\relax\detokenize{#1}\relax\else
      \vspace{\espaciocopiadest\baselineskip}Con copia a:\par
      {\despachotamcopia\despachointercopia #1\par}%
    \fi
    \underline{Presente.}~--\par
  \end{minipage}%
  \par\vspace{\espaciopartes\baselineskip}%
}

%% --- ASUNTO y REFERENCIA (bloque con sangria, etiqueta en negrita) ---
%%  Sangria FRANCESA (Edison 2026-10-10, formato de la bancada): la segunda linea del
%%  texto empieza justo debajo de la primera letra del texto, no bajo la etiqueta; bajo
%%  «ASUNTO:» queda vacio. La etiqueta ocupa una columna comun (la mas ancha de ASUNTO y
%%  REFERENCIA) para que los dos textos queden alineados. El asunto va como se escribe.
\newlength{\despachoEtiquetaAsunto}
\newlength{\despachoEtiquetaNombre}
%%  El TEXTO empieza en \sangriaasunto (la mitad de la hoja); la etiqueta ocupa la columna
%%  inmediatamente a su izquierda (2026-10-10: antes era la etiqueta la que arrancaba ahi).
%%  Los DOS PUNTOS en una misma columna (Edison 2026-10-18): la palabra va en una caja del
%%  ancho de «REFERENCIA» y el «:» detras de esa caja, asi «ASUNTO     :» y «REFERENCIA:»
%%  coinciden (antes el «:» iba pegado a cada palabra). \despachoEtiqueta lo usan los dos.
\newcommand{\despachoEtiqueta}[1]{%
  \makebox[\despachoEtiquetaAsunto][l]{\makebox[\despachoEtiquetaNombre][l]{\textbf{#1}}:}%
}
\newcommand{\despachoMedirEtiquetas}{%
  \settowidth{\despachoEtiquetaNombre}{\textbf{REFERENCIA}}%
  \settowidth{\despachoEtiquetaAsunto}{\textbf{REFERENCIA}:\ }%
}
\newcommand{\bloqueetiquetado}[2]{%
  \despachoMedirEtiquetas
  \noindent\hspace*{\dimexpr\sangriaasunto-\despachoEtiquetaAsunto\relax}%
  \begin{minipage}{\dimexpr\textwidth-\sangriaasunto+\despachoEtiquetaAsunto\relax}%
    \setlength{\parindent}{0pt}\setlength{\parskip}{0pt}%
    \despachoEtiqueta{#1}%
    \parbox[t]{\dimexpr\linewidth-\despachoEtiquetaAsunto\relax}{\raggedright #2}%
  \end{minipage}%
  \par\vspace{\espaciopartes\baselineskip}%
}
%% El texto del asunto va tal como se escribe (solo la etiqueta «ASUNTO:» en mayúsculas);
%% las mayúsculas sostenidas del 2026-10-10 por la mañana se retiran (Edison, misma tarde).
\newcommand{\asunto}[1]{%
  \setmeta{asunto}{#1}%
  \bloqueetiquetado{ASUNTO}{#1}%
}
%% REFERENCIA: palabra completa (no «REF.»), pedido de Edison 2026-09-14.
%% Una sola referencia va en la misma linea de la etiqueta; varias van con
%% \referencias: la etiqueta sola en su linea y la lista a), b), c) DEBAJO,
%% al mismo margen del bloque (como los oficios en Word del despacho). Cada
%% item lleva sangria francesa para que una referencia larga se envuelva
%% bajo su propio texto y no bajo la letra.
\newcommand{\referencia}[1]{\bloqueetiquetado{REFERENCIA}{#1}}
\newcommand{\referencias}[1]{%
  \despachoMedirEtiquetas
  \noindent\hspace*{\dimexpr\sangriaasunto-\despachoEtiquetaAsunto\relax}%
  \begin{minipage}{\dimexpr\textwidth-\sangriaasunto+\despachoEtiquetaAsunto\relax}%
    \setlength{\parindent}{0pt}\setlength{\parskip}{0pt}%
    \despachoEtiqueta{REFERENCIA}%
    \parbox[t]{\dimexpr\linewidth-\despachoEtiquetaAsunto\relax}{\raggedright #1}%
  \end{minipage}%
  \par\vspace{\espaciopartes\baselineskip}%
}
\newcommand{\itemreferencia}[2]{{\hangindent=1.6em\hangafter=1 #1)~#2\par}}

%% --- Saludo protocolar ------------------------------------------------
%%  Antes del saludo, el ASUNTO/REFERENCIA ya dejo \espaciopartes; si el
%%  documento fijo otro «espacio antes del cuerpo» (\setespacioantescuerpo,
%%  core/lib/partes-comunes.tex), se corrige la diferencia aqui (2026-10-01).
\newcommand{\despachoAjusteAntesSaludo}{%
  \ifx\despachoEspacioCuerpo\empty\else
    \vspace{\dimexpr\despachoEspacioCuerpo\baselineskip-\espaciopartes\baselineskip\relax}%
  \fi
}
\newcommand{\saludo}[1][De mi especial consideración:]{%
  \despachoAjusteAntesSaludo
  \noindent #1\par\vspace{\espaciosaludo\baselineskip}%
}

%% --- Cierre -----------------------------------------------------------
\newcommand{\atentamente}[1][Atentamente,]{%
  \vspace{0.4\baselineskip}%
  \noindent\hspace*{\parindent}#1\par
}

%%  \bloquefirma[espacio-para-firma]  (defecto 2.2cm)
%%  Firma \meta{firmante}/\meta{cargofirmante} si existen (personal del
%%  despacho); si no, el titular (\meta{autor}/\meta{cargo}).

%% --- Pie: adjuntos y copias ------------------------------------------

%% (\bloquefirma, \adjunto y \concopia viven en core/lib/partes-comunes.tex)
